\documentclass[UTF8,zihao=5,fontset=windows,a4paper]{ctexart}
\usepackage[left=20mm,right=20mm,top=24mm,bottom=22mm,headheight=20pt,headsep=5mm,footskip=10mm]{geometry}
\usepackage{fontspec,xcolor,graphicx,tikz,booktabs,tabularx,array,caption,fancyhdr,lastpage,enumitem,needspace}
\usepackage{amsmath}
\usepackage[hidelinks,unicode]{hyperref}
\usetikzlibrary{arrows.meta,positioning,calc,fit,backgrounds}
\setmainfont{Times New Roman}
\setsansfont{Calibri}
\setmonofont{Consolas}
\setCJKmainfont{SimSun}[BoldFont=SimHei]
\setCJKsansfont{Microsoft YaHei}
\definecolor{navy}{HTML}{000000}
\definecolor{blue}{HTML}{000000}
\definecolor{ink}{HTML}{000000}
\definecolor{muted}{HTML}{333333}
\definecolor{rulegray}{HTML}{777777}
\definecolor{pale}{HTML}{F7F7F7}
\color{ink}
\hypersetup{pdftitle={Web Agent 的并行执行与容错调度},pdfauthor={陈思旭，寇之洲，谢树韬},pdfsubject={勇闯天涯：方案与工程过程}}
\ctexset{
 section={format=\heiti\fontsize{12}{17}\selectfont,number=\arabic{section},aftername=\quad,beforeskip=10pt,afterskip=5pt},
 subsection={format=\heiti\fontsize{10.5}{15}\selectfont,number=\thesection.\arabic{subsection},aftername=\quad,beforeskip=8pt,afterskip=4pt}
}
\setlength{\parindent}{2em}
\setlength{\parskip}{1.5pt}
\linespread{1.13}
\clubpenalty=10000
\widowpenalty=10000
\pagestyle{fancy}
\fancyhf{}
\fancyhead[L]{\footnotesize Web Agent 的并行执行与容错调度}
\fancyhead[R]{\footnotesize 勇闯天涯}
\fancyfoot[C]{\footnotesize\thepage}
\renewcommand{\headrulewidth}{0.35pt}
\renewcommand{\footrulewidth}{0pt}
\renewcommand{\headrule}{\hbox to\headwidth{\color{rulegray}\leaders\hrule height \headrulewidth\hfill}}
\captionsetup{font=small,labelfont=bf,textfont=normalfont,labelsep=quad,skip=5pt,hypcap=false}
\DeclareCaptionType{algorithm}[算法]
\setlength{\textfloatsep}{10pt}
\setlength{\intextsep}{9pt}
\setlength{\floatsep}{9pt}
\renewcommand{\topfraction}{0.85}
\renewcommand{\bottomfraction}{0.75}
\renewcommand{\textfraction}{0.12}
\setlength{\tabcolsep}{5pt}
\newcolumntype{Y}{>{\raggedright\arraybackslash}X}
\newcommand{\kw}[1]{\textsf{\color{navy}#1}}
\newcommand{\citepaper}[1]{\textsuperscript{\textcolor{blue}{\cite{#1}}}}
\newcommand{\algline}[2]{\makebox[1.4em][r]{\scriptsize\color{muted}#1}\hspace{0.7em}#2\par}
\tikzset{
 box/.style={draw=rulegray,fill=white,line width=.5pt,align=center,inner sep=5pt,font=\small,minimum height=13mm},
 core/.style={box,draw=blue,fill=pale},
 arr/.style={-{Stealth[length=2mm,width=1.4mm]},draw=blue,line width=.75pt},
 note/.style={font=\scriptsize,text=muted,align=center}
}
\begin{document}
\thispagestyle{plain}
\begin{center}
{\heiti\fontsize{18}{25}\selectfont Web Agent 的并行执行与容错调度}\\[8pt]
{\fontsize{11}{16}\selectfont 陈思旭\quad 寇之洲\quad 谢树韬}\\[3pt]
{\small 勇闯天涯}\\[4pt]
{\footnotesize 代码仓库：\href{https://github.com/xie-st/WebAgentNew}{\texttt{github.com/xie-st/WebAgentNew}}}
\end{center}
\vspace{3pt}
\begingroup
\small
\noindent\textbf{摘要：}针对网页任务中的长程状态维护、执行容错和多页面协作问题，本文基于 EffGen 构建统一的网页智能体执行框架，通过可暂停的生成器接口连接官方评测平台与真实浏览器。在此基础上，设计任务账本、过程核验、局部恢复、人工接管及独立只读任务并行机制，并梳理观察筛选、历史压缩与经验复用等候选优化方法。工程实现采用人工规划、文档约束与多智能体分工相结合的流程。
\par\vspace{3pt}
\noindent\textbf{关键词：}网页智能体；并发调度；容错恢复；结果核验；上下文管理
\par
\endgroup
\vspace{4pt}
\section{问题定义与基础工具}
本项目面向网页任务中的并行执行、结果确认、容错恢复和频率控制。在比赛指定的 Qwen3.6-35B-A3B 模型下，我们希望让 Agent 根据页面反馈持续推进任务，在操作失败、结果不明或需要人工介入时保留有效进度，并利用独立页面之间的并行空间缩短执行时间。

开发工具采用 Codex、WorkBuddy 和 DeepSeek Harness（DSH），提供文件读写、代码检索、命令执行、测试及任务派发等通用工程能力。Codex 使用 GPT-6.1 Sol，推理强度为 Medium；WorkBuddy 使用 GLM-5.3-flash；DSH 通过硅基流动 API 调用 DeepSeek v4 flash 1013。上述开发工具不包括本项目所需的 Web Agent 具体知识；赛题理解、论文方法选择和路线规划由团队完成。

运行系统以 EffGen 为基础，复用其观察—思考—工具调用—反馈循环，以及轨迹、预算和上下文压缩机制。\citepaper{effgen} Playwright 提供真实浏览器的页面操作、观察与截图，官方 SDK 提供评测平台的观察、动作和模型接口。我们补充统一网页工具、环境适配和任务状态管理。开发模型构建工程，提交的 Web Agent 使用比赛指定模型执行网页任务。

\section{从人工规划到工程实现}
开发流程如图\ref{fig:workflow}所示。团队先阅读赛题与接口要求，围绕观察效率、长期任务状态、核验恢复和并发调度查询论文，再形成 RoadMap，确定优先级与三阶段路线。Codex 将 RoadMap 翻译成详细方案、DSH 实施说明、验收要求和测试任务，明确输入输出、模块职责、文件归属与前置依赖。

\begin{figure}[htbp]
\centering
\begin{tikzpicture}[x=1cm,y=1cm]
\node[core,text width=3.0cm] (a) at (0,0) {\textbf{团队}\\[-1pt]\footnotesize 赛题 · 论文 · RoadMap};
\node[box,text width=3.0cm] (b) at (4.2,0) {\textbf{Codex}\\[-1pt]\footnotesize 方案 · 实施 · 验收};
\node[box,text width=3.0cm] (c) at (8.4,0) {\textbf{WorkBuddy}\\[-1pt]\footnotesize 编排与派发};
\node[core,text width=3.0cm] (d) at (12.6,0) {\textbf{DSH}\\[-1pt]\footnotesize 实现与离线验证};
\draw[arr] (a)--(b);
\draw[arr] (b)--(c);
\draw[arr] (c)--(d);
\draw[arr] (d.south)--++(0,-.45)-| (a.south);
\node[note,fill=white,inner sep=2pt] at (6.3,-.93) {主线程集成与实际入口验证\quad → \quad 团队反馈与下一批决策};
\end{tikzpicture}
\caption{人机协作的工程开发流程}\label{fig:workflow}
\end{figure}
方案经人工审定后，由 WorkBuddy 交给 DSH 实现与离线验证；主线程负责集成、真实模型与浏览器验证、官方平台接入，每批结果反馈后再确定后续工作。

开发并行沿用模块边界：输入与观察适配、状态与恢复支持、本地任务与判定可独立开发。各写任务使用独立源码副本，测试前确认导入该副本；公共驱动、官方入口、配置和命令行由一个集成负责人顺序合入。强依赖部分先固定接口，真实模型与官方评测顺序运行。

\section{第一阶段：双环境跑通}
第一阶段解决同一个 Agent 如何在官方平台与真实浏览器中连续执行，接入结构如图\ref{fig:architecture}所示。我们固定 EffGen 源码和官方 SDK 版本，建立模型适配器、统一观察与动作协议、会话管理和证据存储，再连接两套环境。

关键改动是在 EffGen 原有循环中加入可暂停的生成器接口：需要执行网页动作时交出动作并暂停，环境执行后将新观察送回同一生成器。官方端每次 \texttt{step()} 返回一个动作，真实端由本地浏览器执行并回传观察；内部状态与证据工具在本地完成。

\begin{figure}[htbp]
\centering
\begin{tikzpicture}[x=1cm,y=1cm]
\node[core,text width=4.2cm,minimum height=18mm] (agent) at (0,0) {\textbf{EffGen 决策循环}\\[-1pt]\footnotesize Qwen3.6-35B-A3B};
\node[box,text width=3.0cm,minimum height=18mm] (tools) at (5.2,0) {\textbf{统一网页工具}\\[-1pt]\footnotesize 动作、观察与元素语义};
\node[box,text width=3.2cm] (official) at (10.2,.9) {\textbf{官方平台}\\[-1pt]\footnotesize SDK · 单动作交互};
\node[box,text width=3.2cm] (real) at (10.2,-.9) {\textbf{真实浏览器}\\[-1pt]\footnotesize Playwright · 渲染与截图};
\draw[arr] ([yshift=3mm]agent.east)--node[above,note]{交出动作}([yshift=3mm]tools.west);
\draw[arr] ([yshift=-3mm]tools.west)--node[below,note]{接回观察}([yshift=-3mm]agent.east);
\draw[arr] (tools.east)--++(.35,0)|- (official.west);
\draw[arr] (tools.east)--++(.35,0)|- (real.west);
\node[note,draw=rulegray,fill=pale,rounded corners=1pt,text width=10.1cm,minimum height=8mm] (evidence) at (3.8,-2.1) {会话与证据存储：任务 · 观察 · 动作 · 反馈 · 模型调用 · 截图};
\draw[draw=rulegray,line width=.6pt] (agent.south)--(agent.south|-evidence.north);
\draw[draw=rulegray,line width=.6pt] (tools.south)--(tools.south|-evidence.north);
\end{tikzpicture}
\caption{共用决策循环的双环境接入结构}\label{fig:architecture}
\end{figure}
两端共用点击、输入、按键、滚动和导航等工具，执行流程见算法\ref{alg:loop}。每个动作绑定当前页面、观察版本和有效元素编号；页面更新或人工接管后重新生成定位信息。日志为后续核验与恢复保留依据。

\begin{algorithm}[htbp]
\hrule height .6pt
\caption{可暂停的观察—动作循环}\label{alg:loop}
\hrule height .3pt
\vspace{4pt}
\begingroup
\small\setlength{\parindent}{0pt}\setlength{\parskip}{2pt}
\algline{1}{初始化 EffGen、会话和预算，安装首个观察 $o$}
\algline{2}{\kw{while} 任务仍活动且预算可用：}
\algline{3}{\hspace*{1.3em}依据 $o$、任务状态和近期反馈推进同一循环，取得提案 $p$}
\algline{4}{\hspace*{1.3em}\kw{if} $p$ 是内部工具：本地执行，将结果送回循环；\kw{continue}}
\algline{5}{\hspace*{1.3em}\kw{if} $p$ 请求完成：}
\algline{6}{\hspace*{2.6em}核验当前要求；通过则结束，否则将缺口送回循环；\kw{continue}}
\algline{7}{\hspace*{1.3em}校验当前观察与动作目标；失败则反馈并 \kw{continue}}
\algline{8}{\hspace*{1.3em}记录待确认操作；\kw{yield} 网页动作，暂停并等待环境}
\algline{9}{\hspace*{1.3em}$o \leftarrow$ 新观察与动作反馈；安装新版本，记录证据}
\endgroup
\vspace{2pt}
\noindent{\rule{\linewidth}{.4pt}}\par
{\footnotesize 注：第 5—6 行对应第二阶段接入的完成核验；第一阶段先建立动作交出与反馈返回。}
\end{algorithm}

\vspace{7pt}
真实浏览器接入比较了 Playwright 与 browser-use 的浏览器层，选择 Playwright 为默认后端。阶段交付覆盖真实模型调用、本地页面任务、公开网页访问、人工暂停继续，以及官方任务的基本运行链路。

\section{第二阶段：基础能力与测试环境}
\subsection{任务状态、结果确认与恢复}
第二阶段的建设方案围绕长任务建立状态与证据之间的联系。账本保存目标、已知事实、缺失信息、待确认操作和证据引用；EffGen 负责规划与行动，程序负责状态转换、版本检查和持久化。模型只表达业务目标与判断，版本号、分支和证据归属由程序绑定。

结果确认放在保存、提交、子目标结束和任务结束等关键节点。可直接判断的条件由程序检查，需要语义判断的条件交给比赛模型。核验未通过时，将缺口反馈给同一循环继续纠正。恢复时依据证据撤回受影响的完成判断，保留其他已确认进度；有副作用的操作超时后，先核实页面结果再决定是否重试。

人工接管覆盖登录、验证码、二次验证和主动暂停。系统停止新动作并处理在途操作后交出浏览器；用户确认继续时重新观察页面，使旧提案和元素编号失效，再依据新状态推进。

测试采用本地可重置网站、官方论坛与电商调测、公开真实网站三层，分别检查可控故障和接管机制、完整任务结果、实际交互与多页面协作。各能力通过独立开关接入实际运行入口。

\Needspace{8\baselineskip}
\subsection{运行时的并发处理}
真实端首版将独立只读任务分配到不同 Tab，每个分支拥有独立循环、会话和观察，共享写操作串行，调度关系如图\ref{fig:scheduler}所示。每个 Tab 最多一个在途动作，一个生成器由单一所有者推进；浏览器操作由统一所有者调度，不将同步页面对象交给多个线程。

并行从两个分支起步，容量可配置；模型调用并发初始为一，先重叠不同页面的等待。所有分支共用根任务预算、站点访问节奏与退避状态。共享账本更新由协调器串行提交，过期版本重新核对。人工接管触发全任务暂停，恢复前统一刷新相关页面。官方端保持单任务、单动作串行交互。

\begin{figure}[htbp]
\centering
\begin{tikzpicture}[x=1cm,y=1cm]
\node[core,text width=13.8cm,minimum height=9mm] (root) at (6.4,0) {\textbf{根协调器}\quad\small 依赖与就绪队列\quad / \quad 共享账本\quad / \quad 总预算与站点限流};
\node[box,text width=4.15cm,minimum height=16mm] (a) at (1.65,-1.4) {\textbf{分支 A · Tab A}\\[-1pt]\footnotesize 独立循环、会话与观察};
\node[box,text width=4.15cm,minimum height=16mm] (b) at (6.4,-1.4) {\textbf{分支 B · Tab B}\\[-1pt]\footnotesize 独立循环、会话与观察};
\node[box,text width=4.15cm,minimum height=16mm] (w) at (11.15,-1.4) {\textbf{共享写任务}\\[-1pt]\footnotesize 串行队列};
\draw[arr] (root.south-|a.north)--(a.north);
\draw[arr] (root.south)--(b.north);
\draw[arr] (root.south-|w.north)--(w.north);
\node[note,text width=14.5cm,minimum height=8mm] (pause) at (6.4,-2.65) {全局接管：停止新动作 → 处理在途操作 → 用户操作 → 重新观察 → 恢复};
\draw[draw=rulegray,dashed] (a.south)--(a.south|-pause.north);
\draw[draw=rulegray,dashed] (b.south)--(b.south|-pause.north);
\draw[draw=rulegray,dashed] (w.south)--(w.south|-pause.north);
\end{tikzpicture}
\caption{独立只读分支并行、共享写串行的真实端调度方案}\label{fig:scheduler}
\end{figure}

\noindent{\heiti\color{navy}开发过程中的依赖与集成}\par
这一分工已用于输入适配、状态恢复和本地判定三线开发。之后，恢复接线、上下文与经验加载、并行调度骨架分开准备；调度的全局暂停仍依赖恢复与接管接口，待接口合入后集成验证。开发并行管理源码归属，运行并行管理页面与业务状态。

\subsection{优化方向与论文来源}
候选方法与论文依据见表\ref{tab:methods}。团队从任务失败方式出发，将论文思路转成适合固定模型与浏览器接口的设计。实施时优先完成基础状态、最小核验和同循环纠错，再接恢复与接管；上下文先保护当前观察和必要账本，提供独立开关与批准文本加载。

\begin{table}[htbp]
\caption{候选优化方向及其论文依据}\label{tab:methods}
\centering
\begingroup
\fontsize{9.4}{13.3}\selectfont
\renewcommand{\arraystretch}{1.17}
\noindent\begin{tabularx}{\linewidth}{@{}p{27mm}p{39mm}Y@{}}
\toprule
\textbf{优化方向} & \textbf{论文依据} & \textbf{采用或候选方式}\\
\midrule
状态与停滞诊断 & PoS / BeliefState\citepaper{pos} & 区分事实、未解决要求与待确认结果；依据目标进展和证据变化诊断停滞。\\
结果核验与纠错 & WALT\citepaper{walt}、Reflexion\citepaper{reflexion} & 明确子目标完成条件，核验反馈回同一循环；短反思记录具体修正办法。\\
同轮摘要与经验 & M²\citepaper{m2} & 候选：行动同轮生成短摘要；经验先加载人工批准文本，默认空库。\\
子目标历史折叠 & AgentProg\citepaper{agentprog} & 候选：折叠已完成过程，把后续仍需使用的对象值和证据保存到账本。\\
观察与历史差分 & Read More, Think More\citepaper{readmore} & 保留当前观察；候选：用差分表达旧观察，减少历史重复内容。\\
页面候选筛选 & Prune4Web\citepaper{prune} & 候选：模型给关键词和权重，固定程序评分，保留当前有效元素映射。\\
多页面依赖调度 & Flash-Searcher\citepaper{flash} & 以依赖图组织独立任务，补充页面归属、共享状态约束和统一限流。\\
截图与行动效率 & UI-TARS\citepaper{uitars}、WebClipper\citepaper{clipper} & 截图先留证；视觉决策、轨迹依赖分析和通用动作组合列为后续候选。\\
\bottomrule
\end{tabularx}
\endgroup
\end{table}

\vspace{6pt}
\noindent{\footnotesize\color{muted}摘要、差分、折叠、筛选及动作组合保留为待比较的方法，最终选择在后续章节填写。}

\section{第三阶段：实验验证}
\vspace{5mm}
\section{实验结果}
\vspace{5mm}
\section{最终设计}
\vspace{5mm}

\vspace{1pt}
\begingroup
\renewcommand{\refname}{参考文献}
\ctexset{section={beforeskip=5pt,afterskip=4pt}}
\fontsize{9}{12}\selectfont
\setlength{\parskip}{0pt}
\begin{thebibliography}{11}
\raggedright
\setlength{\itemsep}{1.7pt}
\setlength{\parsep}{0pt}
\bibitem{effgen} \href{https://arxiv.org/abs/2602.00887}{EffGen: Enabling Small Language Models as Capable Autonomous Agents}. arXiv:2602.00887, 2026.
\bibitem{pos} \href{https://arxiv.org/abs/2610.01415}{Beyond Memory: Harnessing Long-Horizon Agents with Explicit Belief States}. arXiv:2610.01415, 2026.
\bibitem{walt} \href{https://arxiv.org/abs/2510.01524}{WALT: Web Agents that Learn Tools}. arXiv:2510.01524, 2025.
\bibitem{reflexion} \href{https://arxiv.org/abs/2303.11366}{Reflexion: Language Agents with Verbal Reinforcement Learning}. arXiv:2303.11366, 2023.
\bibitem{m2} \href{https://arxiv.org/abs/2603.00503}{M²: Dual-Memory Augmentation for Long-Horizon Web Agents via Trajectory Summarization and Insight Retrieval}. arXiv:2603.00503, 2026.
\bibitem{agentprog} \href{https://arxiv.org/abs/2512.10371}{AgentProg: Empowering Long-Horizon GUI Agents with Program-Guided Context Management}. arXiv:2512.10371, 2025.
\bibitem{readmore} \href{https://arxiv.org/abs/2604.01535}{Read More, Think More: Revisiting Observation Reduction for Web Agents}. arXiv:2604.01535, 2026.
\bibitem{prune} \href{https://arxiv.org/abs/2511.21398}{Prune4Web: DOM Tree Pruning Programming for Web Agent}. arXiv:2511.21398, 2025.
\bibitem{flash} \href{https://arxiv.org/abs/2509.25301}{Flash-Searcher: Fast and Effective Web Agents via DAG-Based Parallel Execution}. arXiv:2509.25301, 2025.
\bibitem{uitars} \href{https://arxiv.org/abs/2501.12326}{UI-TARS: Pioneering Automated GUI Interaction with Native Agents}. arXiv:2501.12326, 2025.
\bibitem{clipper} \href{https://arxiv.org/abs/2602.12852}{WebClipper: Efficient Evolution of Web Agents with Graph-based Trajectory Pruning}. arXiv:2602.12852, 2026.
\end{thebibliography}
\endgroup
\end{document}
